%!TEX root = ../Thesis.tex
\chapter{Afterwords}

I briefly include here some afterwords for anyone who may be interested in them. They reflect a little bit on my thoughts on the things I have presented in the thesis, how I have written my thesis and although very brief, attempt to give a bit of a more nuanced view on particularly Bayesian Deep Learning.\\

\noindent Throughout the writing of this thesis I have tried to show a bit of myself and make the writing somewhat personal rather than only writing about my research. In the introductory chapters to my research papers I included some descriptions of how projects got started, with the hope being that it would put the reader into the mindset of me and my colleagues when projects were started, and providing stronger motivations for the projects. Although I know this may be unorthodox, it was my hope that it would make the reading more engaging for the reader - I hope I have succeeded in this. \\

\noindent At the end of my PhD, I must admit that I at times have become somewhat disillusioned with in particular Bayesian Deep Learning - in large part due to the fact that I often found these methods difficult and fidgety to get to work in practice, and I question whether fighting over decimal points on well-established datasets is fruitful for the field. I would like to iterate that I am by no means above this - much of my research did exactly this. I love the Bayesian framework and its promises, but in all honesty I sometimes question how useful Bayesian Deep Learning is in practice - I hope to get a chance to investigate this in my next position in industry. I hope to be proven wrong in my disillusionment. \\

\noindent In relation to these feelings, if I had to give my past PhD self and the Bayesian machine learning community some advice (which I am not sure I am even slightly qualified to do), it would be to follow the words of Wessel Bruinsma, Andrew Y. K. Foong and Richard E. Turner\footnote{\url{https://mlg.eng.cam.ac.uk/blog/2021/03/31/what-keeps-a-bayesian-awake-at-night-part-2.html}}: \textit{..., when we teach the ideas of probabilistic modelling and inference and write boilerplate in our papers, we’d ask that you pause for breath before trotting off the standard justifications. Instead, consider evangelising a little less about how principled the approach is, mention the importance of exploring different modelling assumptions and their consequences, and stress that the quality of approximate inference should be evaluated in a systematic and principled way.} I for a fact will be thinking about these words in my future career.

